% ==================================================================
%  MEASURE TEMPLATE --- copy me, do not edit in place.
%
%  Covers both kinds of measurement in Section~\ref{sec:clustering_methods}:
%    * a dissimilarity measure d(.,.) --- the input a method consumes;
%    * a validity index --- the output a comparison is built from,
%      whether internal, external or relative.
%  Delete the prompts in Parts 3 and 4 that do not apply to your kind.
%
%    cp template/measure_template.tex \
%       documentation/sections/clustering_methods/measures/NN-<name>.tex
%
%  Then: replace <NAME>, replace the label suffix `template' with the
%  measure's short name, and add one \input line under Section 7.1 in
%  cluster_main.tex.
%
%  This file is NOT included in the build.
%
%  Levels: a measure is a \subsubsection under Section 7.1, so the parts
%  below are \paragraph.  Keep their order unchanged in every copy.
%
%  Code counterpart:
%    xxcluster/measures/dissimilarity/<name>.py, or
%    xxcluster/measures/validation/{internal,external,relative}.py.
%  Part 3 must agree with the class attributes declared there ---
%  `is_metric' and `is_symmetric' for a dissimilarity,
%  `higher_is_better', `range_' and `handles_noise' for an index.
%  See CONTRIBUTING.md, Part 2.
% ==================================================================
\subsubsection{<NAME>}
\label{sec:measure:template}

\paragraph{Overview}
\label{sec:measure:template:overview}
\placeholder{What the measure quantifies, in one paragraph, and which
kind it is: a dissimilarity measure, or a validity index --- and if an
index, whether internal, external or relative. State why it is included
here: which methods need it, or which question about a partition it
answers.}

\paragraph{Definition}
\label{sec:measure:template:definition}
\placeholder{The formal definition, using the symbols declared in the
notation table. Cite the original source. Where the measure is defined
only for particular data types or particular cluster representations,
state that here rather than in a later part.}

\paragraph{Properties}
\label{sec:measure:template:properties}
\placeholder{For a dissimilarity measure: which of the three properties
of Definition~\ref{def:math_metric} it satisfies --- identity and
positivity, symmetry, the triangle inequality --- and which it violates.
This is not a formality. A measure that violates the triangle inequality
is admissible under Definition~\ref{def:cluster_method}, but cannot be
used with a method whose correctness depends on it.

For a validity index: the attainable range, the direction (whether high
or low is better), whether it is corrected for chance, and whether it is
defined when observations are labelled as noise.}

\paragraph{Applicability}
\label{sec:measure:template:applicability}
\placeholder{Which data types the measure accepts --- numeric,
categorical, mixed, incomplete --- and which families of
Section~\ref{sec:background:families} it can serve.

For an index, add: whether it applies to a soft partition directly or
only to a defuzzified one, and how observations labelled as noise are
treated. Excluding them scores only the points a method was confident
about, which flatters that method; whichever choice is made, it must be
stated wherever the index is reported.}

\paragraph{Computation and complexity}
\label{sec:measure:template:complexity}
\placeholder{How the measure is computed, its time and space
complexity, and what that means at \projectname{} data volumes. Note in
particular whether it requires the full pairwise matrix, since that is
quadratic in the number of observations and often the binding
constraint.}

\paragraph{Application to \projectname{} data}
\label{sec:measure:template:application}
\placeholder{Implementation, and where the measure is used: which
methods are fitted with it, or which results are reported under it.
Any parameter it carries, with the value chosen and the reason.}

\paragraph{Behaviour}
\label{sec:measure:template:behaviour}
\placeholder{What the measure rewards, demonstrated rather than
asserted --- on our data, on a benchmark with known structure, or on a
synthetic example.

For an index this is essential, because every index encodes a notion of
what a cluster is. An index built on distances to a centroid rewards
compact, isotropic clusters and penalises elongated or irregular ones,
so it does not rank methods of different families neutrally. Where that
applies, say so here and carry it into
Section~\ref{sec:methodology:threats}.}

\paragraph{Limitations}
\label{sec:measure:template:limitations}
\placeholder{Where the measure is uninformative, misleading, or
undefined --- degenerate partitions, unbalanced cluster sizes, high
dimension, or the concentration of distances discussed in
Section~\ref{sec:dimensionality}. Distinguish what was observed from
what is known from the literature.}

\paragraph{Summary}
\label{sec:measure:template:summary}
\placeholder{Two or three sentences: what this measure is good for,
what it must not be used for, and whether it is part of the reported
protocol of Section~\ref{sec:methodology:metrics} or supplementary.}
